
Large language models produce confident but factually incorrect outputs, yet existing detection methods require multiple forward passes or external knowledge bases. This work proposes Cross-Layer Trajectory Instability (CLTI), a single-pass approach that measures how token representations evolve across late transformer layers. On TruthfulQA MC1 with LLaMA-2-7B (layers 24--31), Normalized Trajectory Instability (NTI) achieves AUROC 0.5657 (5-fold CV, all folds > 0.52). Combining NTI with Convergence Monotonicity Index (CMI) and output entropy yields +7.1\% AUROC improvement over an intercept-only baseline (Fisher combined p = 1.15e-05). However, trajectory metrics fail on low-entropy predictions (AUROC 0.5136, 95\% CI [0.4639, 0.5628] includes chance), and Representational Competition Index (RCI) flip patterns occur in >90\% of both correct and incorrect responses, indicating an architectural rather than epistemic phenomenon. These results demonstrate that trajectory features provide complementary signal for uncertain predictions while identifying constraints that limit the method's scope.

